\section*{Diagnostics: testing a likelihood against its simulation}

An approximate likelihood is a misspecified model, and misspecification leaves a classical signature. At the true parameters the score no longer has zero mean, so the estimator the likelihood defines is biased; and the covariance of the score no longer equals the negative expected curvature, so the uncertainty the likelihood reports is wrong. These are the two Bartlett identities, and their failure is the content of the information-matrix-equality diagnostic of \cite{white1982maximum} and of the sandwich covariance of \cite{huber1967behavior,godambe1960optimum}. In the classical setting those identities can only be tested, through a statistic with a known finite-sample defect, because the data-generating law is unknown. Our setting removes that limitation. The forward process is an exact continuous-time Markov chain (CTMC) and we simulate it exactly, so we hold unlimited data drawn from the very law the likelihood claims to describe, at parameters we set. The classical test therefore becomes a direct measurement: the anchor curvature is computed analytically and the score covariance is estimated by Monte Carlo over independent simulated recordings, with no null hypothesis to reject. The apparatus is standard; what is new here is applying it in this domain and reading its output as a measured property of each approximation rather than as a hypothesis to accept or reject.

We read three quantities off the simulation, in increasing strength, and each is read at a parameter vector that has to be named because there are two of them. The first is the simulation truth $\theta_{\mathrm{sim}}$, the point the recordings were generated from, which is known here and never known in an experiment. The second is the pooled optimum $\theta_{\mathrm{pool}}$, the maximum-likelihood fit of the member under test to all the recordings of one cell at once; it is not known in closed form and is located numerically, by Gauss--Newton with Levenberg damping warm-started at the truth, with a near-zero gradient certifying that the point reached is stationary (Methods). Were the likelihood correctly specified the two would agree up to sampling error. They do not, and their difference $\theta_{\mathrm{pool}}-\theta_{\mathrm{sim}}$ is the misspecification bias itself. Everything below is computed at both anchors and every figure states which one it shows, because they answer different questions: at $\theta_{\mathrm{sim}}$ a nonzero mean score is visible at all, whereas at $\theta_{\mathrm{pool}}$ the score vanishes by construction, so the information comparison is not contaminated by a displaced gradient. % src: sections/06_methods.tex "Parameter estimation and the two anchors" (optimiser settings, warm start, near-zero gradient certificate)

The first is the standardized residual. With $\mu_t$ and $\sigma_t^2$ the likelihood's own predictive mean and variance for interval $t$, the residual $r_t=(y_t-\mu_t)/\sigma_t$ must have mean zero, unit variance, and no autocorrelation across intervals if the one-step predictive distribution is correct \emph{at the parameters it is evaluated at} (Figure~\ref{fig:time} draws the second and the third; the first is carried in a supplement, where it separates no member and so returns no verdict). This probes the predictive one interval at a time. It is the cheapest of the three and the only one an experimentalist can run on a real recording, because it needs no ensemble and no known truth; but it does need a fitted point, since at an arbitrary parameter vector the three properties fail for reasons that have nothing to do with whether the approximation is any good. Running it at the fit has a cost of its own. For a Gaussian likelihood the score is a weighted combination of exactly the quantities $r_t$ and $r_t^2-1$, so setting it to zero imposes as many linear constraints on the residual moments as there are free parameters, and a fit that leaves the baseline and the noise level free absorbs most of the first two properties into itself. What no fit absorbs is the temporal structure. Residual whiteness is therefore the part of this check that stays informative at the optimum, and it is not optional: residual autocorrelation is the same signal that the strongest check quantifies below, though it is a conservative reading of it, running $1.6$ to $1.7$ times smaller than the score autocorrelation for the members in which either is large. % src: projects/eLife_2025/figures/paper_both/figure_3.html (row G caption numbers; lag-one residual and score ACF by member: LSE 0.846/0.856, NR 0.510/0.864, R 0.118/0.191, IR -0.010/-0.004). The first residual moment is Figure 3--figure supplement 3 row A, cited from Results so that the three Figure 3 supplements are cited in ascending order.

The second is the score. Writing $\ell(\theta)$ for the total log-likelihood and $S(\theta)=\nabla_\theta\ell(\theta)$ for its gradient, correct specification requires $\mathbb{E}[S(\theta_{\mathrm{sim}})]=0$ at the simulation truth. A nonzero mean score is what biases the estimate, and projected through the parameter covariance it becomes a displacement in parameter units, $b=\mathbf{H}^{-1}\,\mathbb{E}[S]$, the reader-facing statement of how far the estimate is pushed (Figure~\ref{fig:time}). This projection is a first-order approximation, from the expansion of the score at the optimum given below. Unlike the residual, this one cannot be run on an experiment at all, and for two reasons rather than one. It needs a known truth, and it needs an ensemble to take the expectation over. At the only point an experimenter has, their own fit, the score vanishes by construction and the check is empty. The simulator supplies both, and the confidence interval narrows as the number of simulated recordings grows.

The third and strongest is the information itself. Correct specification requires $\mathrm{Cov}[S]=\mathbf{H}$, the information the likelihood reports, the covariance being taken over independent recordings of the same cell. That is the second thing an experiment cannot supply: one recording gives one score, and a covariance cannot be estimated from a single draw. An approximation can pass the residual and score checks and still violate this one by an order of magnitude. The two failures are independent: whether the estimate is biased is governed by how the conductance is modelled over the interval, while whether the reported information agrees with the information delivered is governed by the recursion. We turn all three checks on the whole cost ladder set out in Theory, from least squares on the deterministic mean current (LSE) at the bottom to the boundary-conditioned filter at the top. The range they have to cover is wide. The members that carry no recursion misreport the information grossly, and among the recursive members the interval conditioning leaves distortions small enough that a diagnostic has to be sharp to resolve them, so the same instrument has to remain informative across both. Because the failure is directional, we record it as a matrix,
\[ \mathbf{C}=\mathbf{H}^{-1/2}\,\mathbf{J}\,\mathbf{H}^{-1/2}, \]
with $\mathbf{J}=\mathrm{Cov}(S)$ the covariance of the total score, cross-interval terms included. This definition is exact algebra given $\mathbf{H}$ and $\mathbf{J}$, and under correct specification $\mathbf{J}=\mathbf{H}$ so $\mathbf{C}=\mathbf{I}$ exactly.
% [SIGN VERIFIED 2026-08-04] Direction confirmed against the producing code and left as
% written. legacy/lapack_headers.h:2519 Lapack_PSD_Normalized_Congruence_Matrix computes
% H^(-1/2) J H^(-1/2) with H the FIRST argument, and legacy/distributions.h:470-478 states that
% the reported Fisher information is the H reference of Likelihood_Information_Distortion. So
% C > 1 means the score varies more than the reported information, i.e. the likelihood
% over-states the information the data carry and under-reports the width of its interval.
% Cross-checked: theory/macroir/docs/Likelihood_Information_Distortion/supplement_information_distortion_main.tex:129-130,
% papers/_program/machinery.md section 4, and the NR reading of 13-17 on a member that ignores
% the correlation and must therefore over-state information. No dissenting copy was found.
We fix the direction by convention: $C_{ii}>1$ means the likelihood under-reports the uncertainty in parameter $i$ and is over-confident, $C_{ii}<1$ means it over-reports and is conservative. The eigen-directions of $\mathbf{C}$ say which combinations of parameters are misreported, which a scalar summary cannot recover.

The anchor $\mathbf{H}$ is the likelihood's own Gaussian Fisher information,
\begin{equation}
  \mathbf{H}=\sum_t\mathbb{E}\!\left[
    \frac{(\nabla\mu_t)(\nabla\mu_t)^\top}{\sigma_t^2}
    +\frac{(\nabla\sigma_t^2)(\nabla\sigma_t^2)^\top}{2\sigma_t^4}\right],
  \label{eq:gaussian-fisher}
\end{equation}
formed from the same predictive moments and their parameter derivatives that the score already requires, at no extra cost. For a macroscopic algorithm the likelihood is Gaussian by construction, so this is the model's own Fisher and $\mathbf{C}$ is the information-matrix-equality test carried out in the model's own frame. Two properties follow. The Gaussian Fisher is positive semidefinite by construction, so the diagnostic never has to be rescued from a numerically indefinite anchor. And it is analytic, so it does not inherit the finite-difference noise of a numerical Hessian, which in this regime is severe: the finite-difference Fisher is widely indefinite per replicate, and only pooling across replicates recovers a definite matrix. By design the anchor is the Gaussian Fisher; the numerical Fisher we compute alongside it serves to gauge how faithful that Gaussian anchor is. Comparing the approximation to its own Fisher is not circular. $\mathbf{H}$ is the information the likelihood claims, $\mathbf{J}$ is the information it delivers against data from the true process, the two are computed from different objects, and their disagreement is the definition of misspecification.

The bottom rung of the ladder needs a qualification at the point where the diagnostic is defined, because its predictive variance is not a gating quantity. Least squares fits the deterministic mean current under a single variance held constant over the whole recording, so the information it reports, $\mathbf{H}$ as written above with $\sigma_t^2$ replaced by one constant $\sigma^2$, is the information of a homoscedastic model. The exact process violates that assumption: the variance of a recorded interval rises and falls with the open population over the course of the concentration jump, and it also depends on the channel number, so the constant-variance premise is wrong by construction in every cell we simulate. For LSE the quantity $\mathrm{Cov}(S)$ read against the reported information therefore measures the size of that break. It is not an information-matrix-equality failure of the same kind as a gating-aware likelihood's, where the predictive variance does follow the gating and the discrepancy is attributable to how the gating is approximated within the interval and across intervals. The Results quote LSE's ratio with that reading, and they do not read an LSE ratio and a gating-aware ratio as two values of one comparable quantity.

The same matrix repairs what it measures. The sandwich covariance
\[ \Sigma=\mathbf{H}^{-1}\,\mathbf{J}\,\mathbf{H}^{-1}=\mathbf{H}^{-1/2}\,\mathbf{C}\,\mathbf{H}^{-1/2} \]
is the corrected parameter covariance: it substitutes the information the likelihood actually carries for the information it reported. This step, and the bias vector $b$ above, are first-order approximations, from a Taylor expansion of the score at the optimum with the empirical Hessian replaced by its expectation $-\mathbf{H}$ (see Supplementary Information). We verify the correction directly, without leaning on the expansion, by comparing the sandwich-corrected ellipse against the empirical covariance of the maximum-likelihood estimates (MLEs) over the simulated recordings (Figure~\ref{fig:clouds}).
% PENDING-FIG2: Figure 2 has not been knitted for the merged roster; this cross-reference stands on the members that have clouds.
A diagnostic that also corrects the error bars it flags is a stronger result than one that only scores algorithms.

The distortion splits into two physically distinct parts (Figure~\ref{fig:plane}--figure supplement 1).
% NOTE: the map version (figure_4_supplement_3.Rmd) is not built; the line version carries the same
% decomposition. Renumbered 3 -> 1 on 2026-08-05, when Figure 4's supplement set was cut from eleven
% to the two the prose cites (see the note beside the declarations in 04_results.tex).
The sample distortion $\mathbf{C}_s=\mathbf{H}^{-1/2}\,\mathbf{J}_s\,\mathbf{H}^{-1/2}$, built from the within-interval score covariance $\mathbf{J}_s=\sum_t\mathrm{Cov}(s_t)$ with $s_t=\nabla_\theta\ell_t$ the score of interval $t$, measures per-sample departure from Gaussianity; to the leading corrections it is the third and fourth cumulants of the interval current, the non-Gaussian shape the Gaussian predictive cannot carry. The correlation distortion $\mathbf{R}=\mathbf{J}_s^{-1/2}\,\mathbf{J}\,\mathbf{J}_s^{-1/2}$ measures what the within-interval covariance misses, the cross-interval correlation of the score, which is exactly the local time correlation of the current that \cite{milescu2005maximum} named as the dominant error of non-recursive fitting; $\mathbf{R}=\mathbf{I}$ when the score carries no cross-interval dependence. The figures resolve the same pair in time and one parameter at a time, and the notation follows the matrices. For the parameter under discussion, $F_t$ is the Gaussian Fisher information of interval $t$ alone, so that the corresponding diagonal entry of $\mathbf{H}$ is $\sum_t F_t$; $J_t=\mathrm{Var}(s_t)$ is the variance of that interval's score. Accumulated to interval $T$ they are $F_T=\sum_{t\le T}F_t$ and $J_T=\mathrm{Var}\bigl(\sum_{t\le T}s_t\bigr)$, and it is the second of these, not the sum of the $J_t$, that the reported information has to match. The two statements $J_t=F_t$ and $J_T=F_T$ are therefore not equivalent, and their divergence is the whole of what Figure~\ref{fig:time} measures. These compose exactly. The exact identity is
\[ \mathbf{C}=\mathbf{K}\,\mathbf{R}\,\mathbf{K}^\top,\qquad \mathbf{K}=\mathbf{H}^{-1/2}\,\mathbf{J}_s^{1/2}, \]
which holds as algebra for any $\mathbf{H}$ and $\mathbf{J}_s$, since substituting the definitions collapses the chain. The factor $\mathbf{K}$ satisfies $\mathbf{K}\mathbf{K}^\top=\mathbf{C}_s$, so it is a legitimate square-root factor of the sample distortion, but it is not the symmetric principal square root, and the form $\mathbf{C}=\mathbf{C}_s^{1/2}\,\mathbf{R}\,\mathbf{C}_s^{1/2}$ holds only where $\mathbf{H}$ and $\mathbf{J}_s$ commute. What composes with no assumption at all is the information volume,
\[ \log\det\mathbf{C}=\log\det\mathbf{C}_s+\log\det\mathbf{R}, \]
an exact identity because the determinant is multiplicative. The per-parameter diagonal, the trace, the eigenvalues and the Frobenius norm do not compose across the split; only the volume does, and only the volume is decomposed in the maps.

The correlation distortion has a reader-facing translation. For any additive statistic $X_t$ with sum $L=\sum_t X_t$, the ratio $\kappa=\mathrm{Var}(L)\big/\sum_t\mathrm{Var}(X_t)$ is unity when the terms are independent and larger when they are positively correlated, and it converts to an effective sample size $T_\mathrm{eff}=T/\kappa$: a non-recursive likelihood believes it has $T$ independent intervals when it effectively has $T/\kappa$. This is the number from the diagnostic most worth quoting, and it is exact by definition. When $\mathbf{R}\approx\kappa\mathbf{I}$ the temporal error is isotropic and this scalar captures it; when $\mathbf{R}$ is anisotropic, or when $\mathbf{C}_s\neq\mathbf{I}$, the distortion cannot be reduced to a sample-size correction and the full matrix is needed.

A matrix is not a number, and two scalars summarise it without pretending to replace it. Let
$\lambda_1,\dots,\lambda_d$ be the eigenvalues of $\mathbf{C}$ on the retained subspace. Their
logarithms have a mean $\bar e = d^{-1}\sum_i \log\lambda_i$ and a variance
$s^2 = \mathrm{Var}_i(\log\lambda_i)$, and the two quantities reported throughout are their
exponentials,
\begin{equation}
  m \;=\; e^{\bar e} \;=\; \Bigl(\textstyle\prod_i \lambda_i\Bigr)^{1/d},
  \qquad
  a \;=\; e^{\,s},
  \label{eq:magnitude-anisotropy}
\end{equation}
so that $m$ is the typical factor by which the reported information is wrong, the geometric mean
of the per-direction factors, and $a$ is the multiplicative spread of those factors from direction
to direction. A member is drawn as a point at $m$ with a bar from $m/a$ to $m\,a$, which is why
the axes are logarithmic and why the spread is written as a factor rather than as a plus-or-minus.
Both are one at correct specification. The single scalar that combines them is the
affine-invariant Riemannian distance from the identity,
$D_{\mathrm{AI}} = \bigl(\sum_i \log^2\lambda_i\bigr)^{1/2} = \sqrt{d\,(\bar e^2 + s^2)}$, which
is used where a single number is wanted and is reported nowhere in place of the pair.
In error-bar units take the square root of any of these, since the distortion is a ratio of
informations: a distortion of $1.15$ is a $7\%$ error on a confidence interval, $2$ is $41\%$, and
$4$ turns a nominal $95\%$ interval into about $68\%$.
% ADDED 2026-08-04. The two summaries were named throughout and defined nowhere, which a
% replication audit reported as cannot-implement. Formulas from the emitter's own header,
% src/core/likelihood.cpp:560-563 ("e-bar = (1/d) sum log lambda_i, signed magnitude of distortion
% on log scale; s^2 = Var(log lambda_i), anisotropy; D_AI = sqrt(d(e-bar^2+s^2)) = sqrt(sum
% log^2 lambda_i)"), and the exponentiation is what the figures report, per
% figures/paper_both/figure_6B.Rmd:25-30. The error-bar scale is that file's lines 36-39.

Two conventions govern the edge cases.
The direction is fixed as above, $C_{ii}>1$ being over-confidence and $C_{ii}<1$ conservatism, stated once so every verdict reads the same way. And the anchor can be near-singular where a parameter is unidentifiable, which by design happens once the open population relaxes (Figure~\ref{fig:time}--figure supplements 1 and 2); there we diagonalize $\mathbf{H}$ and retain the eigenmodes whose eigenvalue exceeds a tolerance set relative to the largest, $\lambda_i > \lambda_{\max}\cdot 10^{-10}$, evaluate every inverse power on that retained subspace alone, and leave the quantity undefined, not finite, when the retained subspace is empty. The dimension $d$ of Eq.~\ref{eq:magnitude-anisotropy} is the dimension of that subspace, so a cell where a direction has been dropped is summarised over fewer directions than one where none has. Undefined is not zero, and the maps render those cells in grey to keep the two apart.
